% TOP_PROBLEM: {"id": 7200020, "title": "Virasoro conjecture", "category_id": 5, "category": {"id": 5, "name": "algebraic_geometry", "display_name": "Algebraic Geometry", "description": "Geometric objects defined by polynomial equations.", "slug": "algebraic-geometry", "order_index": 5, "created_at": "2026-07-31T15:26:25.671Z"}, "status": "partially_solved"}
% FIRST_POSTED: 2026-10-01
% ATTEMPT_STATUS: unresolved
% !TeX program = lualatex
% Definition and sources only; this file is not an LLM solution attempt.
\documentclass[11pt]{article}
\usepackage[margin=25mm]{geometry}
\usepackage{fontspec}
\setmainfont{Latin Modern Roman}
\usepackage{amsmath,amssymb,mathtools}
\usepackage{unicode-math}
\setmathfont{Latin Modern Math}
\usepackage{xurl,hyperref}
\hypersetup{colorlinks=true,urlcolor=blue}
\setcounter{secnumdepth}{0}
\setlength{\parindent}{0pt}
\setlength{\parskip}{5pt}
\setlength{\emergencystretch}{3em}
\begin{document}
\section*{\texorpdfstring{Virasoro conjecture}{Virasoro conjecture}}\label{7200020}
\subsection{Definitions and mathematical statement}
Choose a homogeneous cohomology basis of $X$. Descendant Gromov--Witten correlators integrate products of evaluation classes and powers $\psi_i^{a_i}$ of cotangent-line classes over the virtual fundamental classes of stable-map moduli spaces $\overline{\mathcal M}_{g,n}(X,\beta)$. Their generating functions $F_g$ form
\[
 Z_X=\exp\left(\sum_{g\ge0}\hbar^{g-1}F_g\right)
\]
as a formal series over the Novikov ring recording effective curve classes. The conjecture asserts $L_k^XZ_X=0$ for every $k\ge-1$, where $L_k^X$ is exactly the Eguchi--Hori--Xiong--Katz operator in Getzler, Section 1.4, including constant terms and signs for odd cohomology. Those operator formulas are incorporated by reference rather than replaced by the Virasoro commutation relations alone. Equality is coefficientwise in all descendant variables, genera, and curve classes.
\par\smallskip\textit{Formulation and concept references:} \href{https://arxiv.org/pdf/math/9812026}{[S1]}.

\subsection{Short English statement}
Do the total descendant Gromov–Witten invariants of every smooth projective complex variety satisfy all of the specified Virasoro differential constraints?
\subsection{Research attempt: two operator constraints and a coefficient-level obstruction}
\textbf{Status and source dependence.} The target fixes the actual Eguchi--Hori--Xiong--Katz operators, including odd-cohomology signs. Getzler [S1] proves their Virasoro commutation relations and explicitly records the reduction to $L_{-1}$ and $L_2$. We rederive that reduction and the even-variable coefficient equations to isolate the missing enumerative identity. Neither is asserted as a new theorem about Gromov--Witten invariants.

\subsubsection{1. Why two specified constraints suffice}
For the operators in the source, $[L_m,L_n]=(m-n)L_{m+n}$ in the permitted range. If a formal vector $Z$ obeys $L_{-1}Z=L_2Z=0$, then
\[L_1Z=-\tfrac13[L_{-1},L_2]Z=0,\qquad L_0Z=-\tfrac12[L_{-1},L_1]Z=0.\]
For $k\ge2$, $[L_1,L_k]=(1-k)L_{k+1}$, whose coefficient is nonzero over characteristic zero. Induction gives all higher constraints. These manipulations are coefficientwise compositions in the formal-series convention where the source defines the operators. They do not introduce a convergence claim for an analytic partition function.

The logical role of this computation is limited: the commutators propagate annihilation once it has been proved. They do not establish even the first annihilation. For example, $D_m=-x^{m+1}\partial_x$ on Laurent polynomials satisfies the same commutator formula, but $D_m(x)=-x^{m+1}\ne0$. Thus replacing the target's concrete operators and correlators by an abstract Virasoro representation would lose the problem.

\subsubsection{2. An explicit coefficient computation}
To display the structure without suppressing superalgebra signs, restrict this computation to commuting variables. Let
\[\mathcal L=\frac{\hbar}{2}\sum_{i,j}A_{ij}\partial_i\partial_j+
\sum_i B_i(t)\partial_i+\frac{C(t)}{2\hbar}+D,\qquad
F=\sum_{g\ge0}\hbar^{g-1}F_g,\quad Z=e^F,\]
with coefficients independent of $\hbar$; take finite sums or a coefficientwise well-defined formal operator. The product rule gives the exact identity
\[Z^{-1}\mathcal L Z=\frac{\hbar}{2}\sum_{i,j}A_{ij}(F_{ij}+F_iF_j)+\sum_i B_iF_i+
\frac{C}{2\hbar}+D.\]
Writing $F_{-1}=0$, its coefficient at $\hbar^{g-1}$ is
\[\frac12\sum_{i,j}A_{ij}\left(\partial_i\partial_jF_{g-1}+
\sum_{h=0}^{g}\partial_iF_h\,\partial_jF_{g-h}\right)
+\sum_i B_i\partial_iF_g+\frac12\delta_{g0}C+\delta_{g1}D.\tag{1}\]
In particular, the genus-zero quadratic term survives: it comes from the product of two $\hbar^{-1}$ terms followed by multiplication by $\hbar$. Dropping this contribution, or placing the constant $D$ in genus zero, gives the wrong recursion.

\subsubsection{3. What a recursive proof would still need}
Even when the lower-genus $F_h$ are known, (1) contains terms involving both $F_0$ and $F_g$. It is an equation for $F_g$ with coefficients determined by genus zero, not a statement that an arbitrary $F_g$ satisfies it. A proof of the target must establish those equations for the actual virtual-cycle integrals, including each Novikov degree and descendant monomial. The string equation addresses $L_{-1}$ in the source's conventions; the unreduced task is the genuine $L_2$ identity. The commutator reduction provides no substitute for it.

For odd cohomology, ordered superderivatives and Koszul signs must be restored exactly as in [S1]. The displayed commuting-variable calculation is a diagnostic subcase, not an alternative definition of $L_k^X$. It also does not assume that arbitrary curve classes can be reordered across an infinite sum without formal finiteness.

\subsubsection{4. Verification and gap}
The numerical factors in the first step follow from $-1-2=-3$ and $-1-1=-2$; the induction never divides by the vanishing coefficient at $k=1$. Formula (1) follows by two explicit differentiations and tracking powers of $\hbar$. The remaining gap is to prove $L_2^XZ_X=0$ for every target variety, in the full signed operator convention. No new general constraint is proved here.

\subsection{Sources}
\begin{itemize}
\item[S1]E. Getzler, The Virasoro conjecture for Gromov-Witten invariants, arXiv version 4 (1999).
\url{https://arxiv.org/pdf/math/9812026}.
\textit{Formulation source.}
\end{itemize}
\end{document}
